\documentclass[11pt]{article}
% Le préambule commun charge babel avec [french] seul ; le résumé anglais (otherlanguage{english})
% a besoin de la langue english chargée. On l'ajoute ici sans modifier le préambule commun.
\PassOptionsToPackage{english}{babel}
% Préambule commun des prétirages v2 (format arXiv : source TeX autonome, pdfLaTeX).
% Moteur : pdfLaTeX (arXiv) → fontenc T1 + inputenc ; XeTeX/LuaTeX (tectonic) → fontspec. Sans ce choix, les
% guillemets « » tapés en dur sortent en « ń ż » sous tectonic (ledger, 2026-11-14).
\usepackage{iftex}
\ifPDFTeX
  \usepackage[T1]{fontenc}
  \usepackage[utf8]{inputenc}
\else
  \usepackage{fontspec}
\fi
\usepackage[french]{babel}
\usepackage{lmodern,microtype}
\usepackage[a4paper,margin=2.5cm]{geometry}
\usepackage{amsmath,amssymb}
\usepackage{graphicx,booktabs,siunitx}
\usepackage{authblk}
\usepackage[numbers,sort&compress]{natbib}
\usepackage[hidelinks]{hyperref}
\usepackage{csquotes}
\sisetup{locale=FR,group-minimum-digits=4}
\renewcommand\Affilfont{\small}
% Étiquette de provenance obligatoire sur chaque nombre (règle du projet) :
\newcommand{\mesure}{\textsuperscript{\textsc{m}}}   % mesuré
\newcommand{\derive}{\textsuperscript{\textsc{d}}}   % dérivé
\newcommand{\suppose}{\textsuperscript{\textsc{s}}}  % supposé


\title{Un jumeau numérique de l'offre de soins à contrat vérifiable :\\
le stock-flux reproduit à l'écart nul, et des critères d'acceptation qui savent échouer}
\author[1]{Équipe hackathon-santé}
\affil[1]{Projet « jumeau numérique de l'offre de soins », dépôt \texttt{hackathon-sante/v2}}
\date{Prétirage, version du \today{} --- non relu par des pairs}

\begin{document}
\maketitle

\begin{abstract}
Projeter l'effectif des professions de santé « libérales » (secteur privé) par territoire jusqu'en 2050
suppose un socle commun. Méthodes, paramètres, données et protocoles de validation doivent pouvoir s'y
brancher sans se contredire. Nous décrivons ce socle, un jumeau numérique fait de trois objets :
\begin{itemize}
  \item un état $N[\text{profession},\text{secteur},\text{territoire},\text{sexe},\text{âge}]$, avec un
        vivier en formation ;
  \item des modules de flux en greffons, qui déclarent leurs paramètres sourcés et leurs dépendances ;
  \item des scénarios validés et sérialisables, faits d'hyperparamètres et d'événements datés.
\end{itemize}
Le modèle stock-flux existant, réécrit en trois modules, reproduit l'implémentation de référence avec
un écart de $0{,}0$\mesure{} personne sur 140\mesure{} comparaisons et sur
\num{355248}\mesure{} valeurs cellule $\times$ âge $\times$ année jusqu'en 2050. Une première version
des critères d'acceptation passait par construction. Deux revues adversariales successives l'ont
montré, avec 6\mesure{} puis 4\mesure{} défauts bloquants confirmés en exécutant le code. Nous avons
donc réécrit chaque critère pour qu'il puisse échouer, et vérifié qu'il échoue quand on réintroduit
le défaut (mutation). Une troisième revue ne trouve plus de défaut bloquant ; ses défauts importants
sont corrigés à leur tour. Les 13\mesure{} critères passent ; un scénario se recalcule en
22 à 36\mesure{}~ms au département. La principale limite est le passage à la commune : environ
6,8\derive{}~s extrapolés pour cinq professions et deux secteurs, au-dessus de la cible de
2\suppose{}~s.
\end{abstract}
\begin{otherlanguage}{english}
\begin{abstract}
Projecting France's private-sector health workforce by territory up to 2050 needs a shared substrate
onto which methods, parameters, data and validation protocols can plug without contradicting one
another. We describe that substrate, a digital twin built from three objects:
\begin{itemize}
  \item a state $N[\text{profession},\text{sector},\text{territory},\text{sex},\text{age}]$ plus a
        training pipeline;
  \item plug-in flow modules that declare sourced parameters and their dependencies;
  \item validated, serialisable scenarios made of hyperparameters and dated events.
\end{itemize}
Rewritten as three modules, the existing stock-flow model reproduces the reference implementation
with a gap of $0.0$\mesure{} persons over 140\mesure{} comparisons and \num{355248}\mesure{}
cell~$\times$~age~$\times$~year values to 2050. A first set of acceptance criteria passed by
construction. Two successive adversarial reviews showed this, confirming 6\mesure{} and then
4\mesure{} blocking defects by running the code. We therefore rewrote each criterion so that it can
fail, and checked that it does fail when the defect is reinserted (mutation). A third review found no
blocking defect left; its important findings were fixed in turn. All 13\mesure{} criteria pass, and a
scenario recomputes in 22--36\mesure{}~ms at département level. The main limitation is the move to
commune level: about 6.8\derive{}~s extrapolated for five professions and two sectors, above the
2\suppose{}~s target.
\end{abstract}
\end{otherlanguage}
\noindent\textbf{Mots-clés :} démographie des professions de santé ; projection ; jumeau numérique ;
modèle stock-flux ; architecture en greffons ; tests de mutation ; revue adversariale.

\noindent\small Exposants de provenance : \mesure{} mesuré, \derive{} dérivé, \suppose{} supposé.\normalsize

% ---------------------------------------------------------------------------------------------------
\section{Introduction}

Le projet vise à simuler des politiques publiques nationales, chacune exprimée comme un réglage
d'hyperparamètres sourcés. On mesure leur effet sur l'effectif « libéral » de cinq professions :
sages-femmes, chirurgiens-dentistes, pharmaciens, masseurs-kinésithérapeutes et infirmiers. Les
échelles vont de la commune au national, et l'horizon de référence est 2050. « Libéral » désigne ici
le secteur privé, lucratif ou non : tout employeur public est exclu. Le travail est découpé en six axes
menés en parallèle :
\begin{itemize}
  \item paramètres (axe~1) ;
  \item méthodes (axe~2) ;
  \item validation (axe~3) ;
  \item données (axe~4) ;
  \item sensibilité (axe~5).
\end{itemize}
Sans un objet commun, chaque axe recoderait la dynamique à sa façon, et les résultats ne seraient
pas comparables.

L'axe~0, objet de ce prétirage, fournit cet objet commun. La question est double :
\begin{enumerate}
  \item un contrat de modules peut-il porter le modèle existant \emph{à l'identique}, tout en
        acceptant d'autres méthodes sans modifier aucun fichier ?
  \item comment s'assurer que les critères d'acceptation d'un tel contrat prouvent ce qu'ils
        annoncent ?
\end{enumerate}

\paragraph{Contributions.}
\begin{itemize}
  \item Un contrat à trois objets (état, modules, scénario). Les flux sont typés, les dépendances
        déclarées et vérifiées, et le registre des flux est tenu par cellule (section~\ref{sec:methode}).
  \item La réécriture du modèle stock-flux de référence en modules, avec un écart \emph{nul} à
        l'implémentation d'origine. La v1 sert d'oracle figé (section~\ref{sec:resultats}).
  \item Treize critères d'acceptation écrits pour pouvoir échouer. Chacun est accompagné de la valeur
        mesurée sur le code fautif, et quatre d'entre eux d'une mutation rejouée.
  \item Le récit de trois revues adversariales : ce qu'elles ont trouvé, et ce que cela dit des tests
        « vrais par construction » (section~\ref{sec:validation}).
\end{itemize}

% ---------------------------------------------------------------------------------------------------
\section{Travaux existants}

\paragraph{Projections de la DREES.} Les Dossiers de la DREES n\textsuperscript{o}~76
\citep{drees2021dd76} projettent jusqu'en 2050 les médecins, chirurgiens-dentistes, sages-femmes et
pharmaciens. Le modèle procède par composantes :
\begin{itemize}
  \item il est ventilé par âge, sexe et mode d'exercice ;
  \item il représente les cessations, les transitions de mode d'exercice, les reprises, les entrées
        décalées de la durée des études et les diplômés étrangers.
\end{itemize}
Nous reprenons la même décomposition en flux. Notre contrat n'impose cependant pas une méthode : il
fixe une interface que plusieurs méthodes peuvent remplir. La reproduction du tendanciel DD76 en 2050
est un critère de validation de l'axe~3, pas de l'axe~0.

\paragraph{Implémentation de référence (v1).} Le projet disposait déjà d'un modèle stock-flux par
composantes, sur des tranches d'âge quinquennales, ajusté par profession $\times$ département $\times$
sexe. Ses taux de rétention nette sont rapprochés du niveau national quand les effectifs sont faibles
(poids a priori de 200\suppose{} personnes). Évalué au backtest, il obtient une erreur absolue moyenne
en pourcentage (MAPE) de 8,08\mesure{}~\% par département, contre 8,68\mesure{}~\% pour une tendance
linéaire, sur l'ancien périmètre « libéral + mixte ». Cette v1 est figée, et sert d'oracle.

\paragraph{Observatoire de l'équipe.} Un observatoire développé en parallèle par un membre de
l'équipe\footnote{Dépôt \texttt{Pierreafeu181/Hackathon-Evolution-soin}, commit du 2026-11-14, lu par
le projet ; voir le rapport interne \texttt{reports/2026-11-14-reprise-observatoire.md}.} fournit la
forme corrigée du levier retraite, que nous adoptons (B1, section~\ref{sec:methode}).

\paragraph{Ce que nous n'avons pas cité.} L'idée de réintroduire volontairement un défaut pour
vérifier qu'un test le détecte relève des tests de mutation. Nous n'avons pas pu lire la publication
fondatrice de ce domaine dans le temps imparti : nous ne la citons donc pas.

% ---------------------------------------------------------------------------------------------------
\section{Données}

\begin{table}[h]
\centering\small
\begin{tabular}{p{4.2cm}p{3.3cm}p{6.4cm}}
\toprule
Jeu & Producteur, millésime & Usage et contrôles \\
\midrule
RPPS, dép.\ $\times$ âge $\times$ sexe $\times$ mode, 2012--2026 & DREES, publication 2026-07-02 &
Historique d'ajustement ; oracle v1 via \texttt{sante.model.stock\_table}. \\
Secteur d'activité public / privé & axe~4, d'après DREES et FINESS (extraction 2026-05-04) &
\texttt{history.sector\_table} : les salariés « ambigus » sont répartis par une part privée
centrale, minimale ou maximale. \\
CartoSanté par commune, 2025 & CartoSanté (Atlasanté) &
Test d'agrégation : \num{13147}\mesure{} communes, 3 professions, \num{224339}\mesure{} libéraux. \\
Inscrits par année d'étude, 2020--2025 & DREES (990\_drees\_mmop2) &
Vivier en formation (dentistes, pharmaciens). \\
\bottomrule
\end{tabular}
\caption{Données lues par le jumeau. Toutes sont ouvertes, et chacune est citée avec son millésime.}
\end{table}

\paragraph{Écart externe mesuré.} En 2025, la somme des communes CartoSanté est inférieure au RPPS
départemental (libéral et mixte) :
\begin{itemize}
  \item de 2,8\mesure{}~\% pour les dentistes ;
  \item de 5,4\mesure{}~\% pour les kinés, dont $-$\num{2084}\mesure{} dans le département~13, où
        Marseille manque.
\end{itemize}
La MAPE par département vaut 4,2\mesure{}~\% et 4,7\mesure{}~\%. Cet écart est publié, pas corrigé.

% ---------------------------------------------------------------------------------------------------
\section{Méthode}\label{sec:methode}

\subsection{L'état}
L'état est un tableau $N_{c,a}(t)$. La cellule $c$ est un quadruplet (profession, secteur,
territoire, sexe), et $a$ une tranche d'âge.
\begin{itemize}
  \item \textbf{Le secteur est un axe de l'état, pas un filtre.} Un entrant bloqué, ou une transition
        privé $\leftrightarrow$ public, change de cellule au lieu de disparaître.
  \item \textbf{Le territoire a un niveau déclaré} : commune, département, région ou national. Une
        table de passage agrège vers le haut et refuse toute unité non rattachée, ou rattachée deux
        fois.
  \item \textbf{Le vivier en formation est national} : profession $\times$ année d'étude.
\end{itemize}

\subsection{Les modules}
Chaque module déclare :
\begin{itemize}
  \item un nom et un rang d'application ;
  \item des paramètres, chacun avec valeur, unité, source, millésime, statut (mesuré, sourcé, supposé
        ou non mesurable) et bornes ;
  \item les flux de l'année qu'il lit (\texttt{needs}) et ceux qu'il produit (\texttt{produces}).
\end{itemize}
Il ne modifie jamais l'état. Il rend des flux typés :
\begin{itemize}
  \item \textbf{Shift} (vieillissement) :
        $N_a \leftarrow (1-f_a)N_a + f_{a-1}N_{a-1}$, et la tranche ouverte garde tout ;
  \item \textbf{Rate} : $N \leftarrow s \odot N$, où $(1-s)\odot N$ va dans un compte de sortie ;
  \item \textbf{Add} : $N \leftarrow N + I$ ;
  \item \textbf{Move} : transfert d'une cellule à une autre ;
  \item \textbf{Pool} : mise à jour du vivier, avec des comptes pour les diplômés et les abandons.
\end{itemize}
Le moteur applique les flux module par module : c'est une séparation d'opérateurs. Il applique
ensuite un plancher à 0, et inscrit chaque flux par cellule dans un registre.

À la construction, il vérifie que chaque besoin est produit par un module placé avant, et qu'aucun
flux n'a deux producteurs. Pendant le run, il refuse toute lecture et tout flux non déclarés. Ajouter
un module revient à déposer un fichier : aucun autre fichier n'est modifié. Une méthode qui n'est pas
un stock-flux s'y branche de la même façon, en rendant ses résultats sous forme de flux.

\paragraph{Stock-flux de référence en trois modules.} Notons $b = \text{Shift}(N)$, avec $f = 1/5$
pour toutes les tranches (convention v1). On a alors
\begin{equation}
N(t+1) = \max\!\bigl(0,\; s\odot b + k(t)\,(I + p\,h)\bigr),
\end{equation}
où :
\begin{itemize}
  \item $s \le 1$ est la rétention nette ajustée ;
  \item $I$ est l'entrée nette annuelle et $p$ sa pente ;
  \item $k(t)$ est le multiplicateur de filière ;
  \item $h = t - t_0$ est l'horizon.
\end{itemize}
Les opérations flottantes sont celles de la v1, dans le même ordre, et l'ajustement réutilise sa
fonction \texttt{fit\_rates}. L'écart nul est donc \emph{construit} : il prouve que le contrat ne
déforme rien, pas que le modèle est juste.

\paragraph{Levier retraite (B1).} Un report de $\Delta$ ans modifie le taux de sortie $e = 1-s$ :
\begin{itemize}
  \item tranches de 65~ans et plus : $e' = e/(1+\lambda e)$. Avec une seule tranche ouverte,
        $\lambda = \Delta$, et le séjour moyen passe de $1/e$ à $1/e + \Delta$. Avec plusieurs
        tranches (65--69, 70--74, 75+), $\lambda$ est calculé par cellule. On le choisit pour que le
        temps espéré en 65+ d'une personne qui entre en 65--69 augmente de $\Delta$ exactement. Ce
        temps est l'espérance d'une chaîne de Markov dont les transitions sont le vieillissement $f$ et
        la rétention $s$, dans l'ordre du moteur ;
  \item 60--64~ans : $e' = e(1-\Delta/5)$ ;
  \item 55--59~ans : $e' = e(1-\Delta/10)$ ;
  \item tranches à solde net positif ($e = 0$) : inchangées.
\end{itemize}
La montée en charge est de 3\suppose{}~mois par an, à partir de l'année qui suit la décision.

\paragraph{Régulation de l'installation.} Dans une zone en « remplacement seul », les entrées
privées (positives) d'une année ne peuvent dépasser les sorties de cette année.
\begin{itemize}
  \item \textbf{Fuite.} Seule une part $\theta = 0{,}80$\suppose{} de l'excédent est bloquée (axe~3,
        expérience naturelle sur le zonage infirmier, exploratoire).
  \item \textbf{Destination des bloqués.} Elle se répartit entre trois issues :
        \begin{itemize}
          \item un territoire voisin non régulé, dans le privé (défaut : 100\suppose{}~\%) ;
          \item le salariat, sous la forme d'un secteur \texttt{bloques} qui suit la démographie du
                privé par une règle de repli explicite ;
          \item la sortie, inscrite dans un compte du registre.
        \end{itemize}
  \item \textbf{Effet au département.} La part « voisine » y reste dans sa cellule : la régulation par
        défaut est donc sans effet à ce grain. À la commune, elle exige une table de voisinage.
  \item \textbf{Exemption.} Les dentistes spécialistes sont exemptés, mais leur part n'est pas
        mesurée : elle vaut 0 par défaut.
\end{itemize}
La destination réelle des entrants bloqués n'est chiffrée par aucune source. L'option « vers le
public » est publiée comme mesure de sensibilité.

\subsection{Le scénario}
Un scénario contient :
\begin{itemize}
  \item des surcharges \texttt{module.param} ;
  \item des événements datés (paramètre, début, fin, opération \texttt{set}, \texttt{mul} ou
        \texttt{add}, sélecteur sur les axes de l'état) ;
  \item une composition de modules et une graine.
\end{itemize}
Il est \emph{posé sur} le scénario de base de l'estimateur : ses paramètres priment, ses événements
s'ajoutent. Avant le run, il est validé. Sont refusés :
\begin{itemize}
  \item les noms de paramètres, opérations et axes inconnus ;
  \item les valeurs de sélecteur absentes de l'état ;
  \item les valeurs hors des bornes déclarées.
\end{itemize}
Il se sérialise en JSON, avec un numéro de format. Les quotas suivent la convention v1 : un
multiplicateur décidé à $t_0$ agit à partir de $t_0 + L_p + 1$, où $L_p$ est la durée des études.

% ---------------------------------------------------------------------------------------------------
\section{Résultats}\label{sec:resultats}

\begin{table}[h]
\centering\small
\begin{tabular}{lp{6.3cm}p{5.0cm}}
\toprule
\# & Critère et mesure après correction & Valeur sur le code fautif \\
\midrule
C1 & Reproduction v1 : écart $0{,}0$\mesure{} sur 140 comparaisons de \texttt{predict} et sur
\num{355248} valeurs (7 professions, jusqu'en 2050) & non contesté \\
C2 & Régulation (état réel, 2 secteurs, 100~\% salariat, $\theta = 0{,}8$) : privé des zones
$-$\num{6709}\mesure{} (kinés) et $-$\num{3470}\mesure{} (dentistes) ; $\Delta$ total sans plancher
$0{,}0$\mesure{} & effet nul ; puis $\Delta$ total $+254$ et $-474$ \\
C3 & Module factice déposé hors du paquet, exécuté : flux \num{4040} = attendu & test hors de la
voie du banc \\
C4 & Régulation sur le privé seul : $\Delta$ total $-2{,}9\cdot10^{-11}$\mesure{} & $+822$ kinés \\
C5 & Dépendance manquante, double producteur, lecture non déclarée, Pool lu par erreur : 4 refus
explicites & résultat faux sans erreur \\
C6 & Faute de frappe, sélecteur \texttt{'dentiste'} ou \texttt{6}, quotas $\times110$ : refusés ;
aller-retour JSON : écart $0{,}0$ & ignorés ; $\times110 \to +$\num{5085033} \\
C7 & Quotas $\times1{,}1$ : écart $4{,}5\cdot10^{-13}$\mesure{} avec la v1. Retraite : séjour en
65+ $+1{,}000 / +2{,}000 / +5{,}000$\mesure{}~an pour $\Delta = 1/2/5$ & quotas $+226$ ; séjour
$+1{,}31$ an, infini pour $\Delta \ge 5$ \\
C8 & Agrégation commune $\to$ national : écart interne 0\mesure{} ; unité inconnue refusée &
comparaison circulaire \\
C9 & Aucun état partagé entre deux runs (vivier $\times2 \to$ diplômés $\times2$) & état
persistant \\
C10 & Vieillissement à largeurs variables : total conservé & $1/5$ codé en dur \\
C11 & Scénario posé sur la base (variante filière) : quotas $\times1{,}1 \to +6{,}53$\mesure{}~\% &
$-15{,}2$~\% \\
C12 & Paramètres sourcés : $\theta = 0{,}8 \to$ 0,800\mesure{} des bloqués ; exemption 50~\%
$\to$ 0,500\mesure{} ; défaut 100~\% voisine au département $\to$ écart 0 ; abandon IDE
0,105\derive{} ; cumul par profession & abandon 0 ; cumul 0,10 partout \\
C13 & Bornes sur la valeur composée et sur les paramètres ajustés ; JSON + \texttt{over} ;
règle « entrées $\le$ sorties » : écart $7\cdot10^{-15}$\mesure{} à une formule indépendante ;
retraite à 3 tranches $\ge 65$ : écart de séjour $10^{-14}$\mesure{} & $\times243 \to +14\,317$~\% ;
règle non testée ; séjour $\ne +\Delta$ \\
\bottomrule
\end{tabular}
\caption{Critères d'acceptation de l'axe~0 (\texttt{jumeau.check}, code de sortie~1 en cas
d'échec). Colonne de droite : valeur mesurée par les revues adversariales sur le code d'avant.}
\label{tab:criteres}
\end{table}

\paragraph{Effets des leviers (dentistes, 2050).}
\begin{itemize}
  \item Retraite +1~an : $+1{,}60$\mesure{}~\% ; +2~ans : $+3{,}22$\mesure{}~\%.
  \item Quotas $\times1{,}1$ décidés en 2026 ($L_p = 6$\suppose{}) : $+6{,}53$\mesure{}~\% sur la
        variante filière.
\end{itemize}
Ces chiffres illustrent le contrat. Ce ne sont pas des résultats de politique publique : les plages
et la sensibilité relèvent de l'axe~5.

\paragraph{Temps de recalcul jusqu'en 2050.}
\begin{itemize}
  \item de 22 à 36\mesure{}~ms pour 7 professions au département (\num{1410} cellules), sur une
        machine partagée avec d'autres sessions ;
  \item à la taille commune (une profession, \num{35047} territoires, données répliquées) :
        0,57\mesure{}~s pour un secteur et 1,36\mesure{}~s pour deux (0,78~s avant les corrections
        de la troisième revue) ;
  \item pour cinq professions et deux secteurs : environ 6,8\derive{}~s par extrapolation, au-dessus
        des 2\suppose{}~s visés.
\end{itemize}

% ---------------------------------------------------------------------------------------------------
\section{Validation et menaces à la validité}\label{sec:validation}

\subsection{Deux revues adversariales}
La première version du jumeau annonçait quatre critères remplis. Une revue adversariale a été menée
par 9 agents en lecture seule, sous trois angles :
\begin{itemize}
  \item les tests peuvent-ils échouer ;
  \item la physique du modèle ;
  \item le contrat comme socle.
\end{itemize}
Les deux constats les plus graves de chaque angle ont été soumis à un agent chargé de les réfuter
en exécutant le code. La revue a confirmé 6\mesure{} défauts sur 6, dont 4 bloquants :
\begin{itemize}
  \item un module ajouté au scénario était ignoré par la voie du banc ;
  \item le secteur public n'avait pas de sorties ($+822$\mesure{} kinés en 2050) ;
  \item les dépendances entre modules étaient implicites ;
  \item le scénario n'était ni validé ni sérialisable.
\end{itemize}
Trois critères étaient vrais par construction. Par exemple, la vérification « aucun fichier
modifié » écrivait le module factice dans un dossier temporaire, si bien qu'elle ne pouvait pas
échouer.

La seconde revue, menée selon la même méthode sur le code corrigé, a confirmé 4\mesure{} nouveaux
défauts bloquants :
\begin{itemize}
  \item le levier retraite touchait la tranche ouverte 65+ en permanence, ce qui rend le séjour
        infini pour $\Delta \ge 5$. Le critère de comparaison à la v1 ne pouvait pas le voir, puisque
        la v1 porte le même défaut ;
  \item la régulation faisait varier le total ;
  \item la correction précédente avait introduit un scénario qui écrasait la base ;
  \item les valeurs de sélecteur n'étaient pas validées.
\end{itemize}

\paragraph{Leçon.} Comparer un mécanisme à une implémentation qui porte le même défaut ne valide
rien. Un critère utile se compare à la grandeur que le mécanisme prétend produire, comme le séjour
attendu, ou à un invariant, comme la conservation des effectifs. Il est accompagné de sa valeur sur
le code fautif.

\subsection{Mutations rejouées}
Pour chaque défaut de la seconde revue, nous avons réintroduit le comportement fautif dans le code
corrigé, puis relancé le critère, qui échoue chaque fois :
\begin{itemize}
  \item ancien masque retraite $\to$ C7 échoue (séjour $+2{,}03$\mesure{}, $+5{,}41$\mesure{},
        infini) ;
  \item destination publique $\to$ C2 échoue ($\Delta$ total $-474$\mesure{} et $+254$\mesure{}) ;
  \item scénario qui remplace la base $\to$ C11 échoue (\num{82374}\mesure{} contre
        \num{102987}\mesure{}) ;
  \item validation des noms seulement $\to$ C6 échoue.
\end{itemize}
Pour la première revue, deux mutations ont été rejouées :
\begin{itemize}
  \item un public sans sorties fait échouer C4 ($+838$\mesure{}) ;
  \item des quotas avancés d'un an font échouer C7.
\end{itemize}

\subsection{Troisième revue}
La troisième revue, menée selon la même méthode, ne trouve plus de défaut bloquant dans le jumeau.
Mesurée sur trois professions, la conservation tient à 0,5\mesure{} personne près en 2050. La revue
confirme quatre défauts importants :
\begin{itemize}
  \item des bornes contournables en empilant des événements ($\times243 \to +14\,317$\mesure{}~\%),
        et NaN accepté ;
  \item des bornes non appliquées aux paramètres ajustés (rétention $> 1$) ;
  \item une composition perdue après un aller-retour JSON ;
  \item un séjour qui ne vaut pas $+\Delta$ quand il y a plusieurs tranches $\ge 65$ (phase~2).
\end{itemize}
Elle relève aussi trois constats non vérifiés. Le plus important est que C2 et C4 ne testaient pas la
règle « entrées $\le$ sorties » elle-même : une mutation qui bloque toutes les entrées les passait.

Tous ces points sont corrigés et couverts par C13. La règle de régulation est désormais comparée à
une formule indépendante, sur une année rejouée à la main. Chaque défaut, réintroduit par mutation,
fait échouer C13 : bornes, JSON, retraite à plusieurs tranches, blocage de toutes les entrées et
déplacement d'entrées négatives.

\subsection{Menaces à la validité}
\begin{itemize}
  \item \textbf{C1 ne valide pas le modèle.} Il prouve seulement que le contrat ne déforme rien. La
        précision reste celle du backtest v1 (8,08\mesure{}~\% de MAPE par département, ancien
        périmètre).
  \item \textbf{La destination \texttt{bloques} est une hypothèse.} Elle conserve le total par
        construction. Si les entrants bloqués rejoignent le public, l'effet total vaut
        $-474$\mesure{} kinés ou $+254$\mesure{} dentistes.
  \item \textbf{La régulation par défaut est sans effet au département} : 100~\% des bloqués vont à
        un territoire voisin, ce qui est supposé. Son effet n'apparaît qu'avec une autre répartition,
        ou à la commune avec une table de voisinage.
  \item \textbf{Le plancher à 0, hérité de la v1, crée des personnes} là où les entrées nettes
        ajustées sont négatives : de 0,02\mesure{} à 0,43\mesure{}~\% de $N(2050)$ selon la
        profession.
  \item \textbf{Le calibrage suppose un vieillissement de 1/5 par an.} Il ne convient pas aux tranches
        communales de largeurs variables.
  \item \textbf{La sensibilité à l'ordre des modules n'est pas mesurée.}
  \item \textbf{Les zones de régulation des tests sont illustratives} (5 départements) : ce ne sont
        pas les zonages conventionnels réels.
\end{itemize}

% ---------------------------------------------------------------------------------------------------
\section{Reproductibilité}
\begin{itemize}
  \item \textbf{Code :} \texttt{v2/jumeau/}. La v1 figée, \texttt{sante/}, sert d'oracle.
  \item \textbf{Commande} (depuis la racine du dépôt) :
        \texttt{PYTHONPATH=v2 uv run python -m jumeau.check}. Elle dure environ une minute et sort
        avec le code~1 en cas d'échec.
  \item \textbf{Sorties :} \texttt{v2/artifacts/a0\_jumeau/results.json}, \texttt{reproduction.csv}
        et \texttt{check.log}.
  \item \textbf{Graine :} chaque scénario porte sa graine (\texttt{ctx.rng}). Les critères de l'axe~0
        sont déterministes.
  \item \textbf{Données :} \texttt{data/processed/} (v1) et \texttt{v2/data/processed/} (axe~4), avec
        les millésimes donnés en section~3.
  \item \textbf{Enregistrements :} rapport \texttt{v2/docs/lab/reports/2026-11-14-jumeau-axe0.md}
        (version~3) et décisions « Contrat du jumeau », « Corrections du contrat » et « Seconde série
        de corrections » dans \texttt{v2/docs/lab/decisions.md}.
\end{itemize}

% ---------------------------------------------------------------------------------------------------
\section{Conclusion}
Le jumeau fournit aux autres axes un contrat vérifiable :
\begin{itemize}
  \item il reproduit exactement le modèle existant ;
  \item il exécute les modules que le scénario déclare ;
  \item il refuse les erreurs au lieu de les avaler ;
  \item il conserve les effectifs lors des déplacements.
\end{itemize}
L'enseignement principal porte sur la validation elle-même. Deux fois, des critères qui passaient
ne prouvaient rien. La parade adoptée a trois volets : écrire chaque critère contre une grandeur
attendue ou un invariant, y joindre sa valeur sur le code fautif, et rejouer la mutation.

\paragraph{Questions ouvertes, avec leur coût} (temps de travail estimé, \suppose) :
\begin{itemize}
  \item troisième revue adversariale : environ 3~min de machine ;
  \item recalcul sur un état communal réel, avec vectorisation si nécessaire : environ 20~min ;
  \item calibrage à fractions par tranche pour les communes : environ 30~min ;
  \item diplômés du vivier reliés aux entrées : environ 1~h ;
  \item plancher converti en sortie explicite : environ 20~min.
\end{itemize}

\bibliographystyle{plainnat}
\bibliography{refs}

\appendix
\section{Annexe : interface d'un module}
\begin{verbatim}
@register
class MonModule(Module):
    name, order = "mon_module", 35
    needs = ("entrees",)            # flux de l'année lus via ctx.delta
    params = (Param("x", 1.0, "unité", "source", "millésime", "supposé",
                    "description", (0.0, 3.0)),)
    def fit(self, hist): ...        # calibrage sur l'historique <= origine
    def step(self, state, ctx):     # lit l'état, ne le modifie pas
        return [Add(self.name, ..., src="mon_compte")]
\end{verbatim}
\end{document}
